\subsection{Data-Driven Mobility Simulation for Digital Twins}
Mobility ``Digital Twins'' motivate data-driven simulators that can reproduce stochastic human behaviors at scale \citep{Wang2022MobilityDigitalTwin}. Traditionally, mobility simulation relied on Agent-Based Models (ABM) and hand-crafted kinematic rules \citep{helbing1995social}, which are interpretable but often hard to calibrate city-wide.

Early data-driven approaches framed trajectory prediction as deterministic regression, using recurrent models or Transformers \citep{alahi2016social,zhou2022hivt,shi2022mtr}. While strong on mean error, deterministic regression collapses multi-modality into a single conditional average. Generative models (VAEs/GANs/diffusion) were introduced to recover diversity \citep{gupta2018socialgan,ho2020ddpm,zhu2023difftraj}, but maintaining \emph{physical validity} remains challenging: samples can violate kinematic realism or exhibit macroscopic shrinkage, reducing their utility for simulation-grade analysis.

\subsection{Physics-Informed Inductive Biases in Urban Computing}
Physics-Informed Machine Learning (PIML) provides inductive biases to help data-driven mobility models respect urban dynamics. Recent diffusion variants incorporate explicit structural constraints via graph priors or road-network-constrained generation \citep{shi2024gdp,wei2024diffrntraj}, and more general priors can be introduced through fields or PDE-inspired guidance \citep{bastek2025pidm}.

Our work targets \emph{macroscopic flow consistency}. We use a train-only navigation field as a lightweight mean-flow prior, and combine it with a residual decomposition so that physical scale is anchored by a deterministic prior while diffusion focuses on stochastic deviations.
